% Options for packages loaded elsewhere
\PassOptionsToPackage{unicode}{hyperref}
\PassOptionsToPackage{hyphens}{url}
\PassOptionsToPackage{dvipsnames,svgnames,x11names}{xcolor}
%
\documentclass[
]{article}

\usepackage{amsmath,amssymb}
\usepackage{iftex}
\ifPDFTeX
  \usepackage[T1]{fontenc}
  \usepackage[utf8]{inputenc}
  \usepackage{textcomp} % provide euro and other symbols
\else % if luatex or xetex
  \usepackage{unicode-math}
  \defaultfontfeatures{Scale=MatchLowercase}
  \defaultfontfeatures[\rmfamily]{Ligatures=TeX,Scale=1}
\fi
\usepackage{lmodern}
\ifPDFTeX\else  
    % xetex/luatex font selection
\fi
% Use upquote if available, for straight quotes in verbatim environments
\IfFileExists{upquote.sty}{\usepackage{upquote}}{}
\IfFileExists{microtype.sty}{% use microtype if available
  \usepackage[]{microtype}
  \UseMicrotypeSet[protrusion]{basicmath} % disable protrusion for tt fonts
}{}
\makeatletter
\@ifundefined{KOMAClassName}{% if non-KOMA class
  \IfFileExists{parskip.sty}{%
    \usepackage{parskip}
  }{% else
    \setlength{\parindent}{0pt}
    \setlength{\parskip}{6pt plus 2pt minus 1pt}}
}{% if KOMA class
  \KOMAoptions{parskip=half}}
\makeatother
\usepackage{xcolor}
\setlength{\emergencystretch}{3em} % prevent overfull lines
\setcounter{secnumdepth}{5}
% Make \paragraph and \subparagraph free-standing
\makeatletter
\ifx\paragraph\undefined\else
  \let\oldparagraph\paragraph
  \renewcommand{\paragraph}{
    \@ifstar
      \xxxParagraphStar
      \xxxParagraphNoStar
  }
  \newcommand{\xxxParagraphStar}[1]{\oldparagraph*{#1}\mbox{}}
  \newcommand{\xxxParagraphNoStar}[1]{\oldparagraph{#1}\mbox{}}
\fi
\ifx\subparagraph\undefined\else
  \let\oldsubparagraph\subparagraph
  \renewcommand{\subparagraph}{
    \@ifstar
      \xxxSubParagraphStar
      \xxxSubParagraphNoStar
  }
  \newcommand{\xxxSubParagraphStar}[1]{\oldsubparagraph*{#1}\mbox{}}
  \newcommand{\xxxSubParagraphNoStar}[1]{\oldsubparagraph{#1}\mbox{}}
\fi
\makeatother


\providecommand{\tightlist}{%
  \setlength{\itemsep}{0pt}\setlength{\parskip}{0pt}}\usepackage{longtable,booktabs,array}
\usepackage{calc} % for calculating minipage widths
% Correct order of tables after \paragraph or \subparagraph
\usepackage{etoolbox}
\makeatletter
\patchcmd\longtable{\par}{\if@noskipsec\mbox{}\fi\par}{}{}
\makeatother
% Allow footnotes in longtable head/foot
\IfFileExists{footnotehyper.sty}{\usepackage{footnotehyper}}{\usepackage{footnote}}
\makesavenoteenv{longtable}
\usepackage{graphicx}
\makeatletter
\def\maxwidth{\ifdim\Gin@nat@width>\linewidth\linewidth\else\Gin@nat@width\fi}
\def\maxheight{\ifdim\Gin@nat@height>\textheight\textheight\else\Gin@nat@height\fi}
\makeatother
% Scale images if necessary, so that they will not overflow the page
% margins by default, and it is still possible to overwrite the defaults
% using explicit options in \includegraphics[width, height, ...]{}
\setkeys{Gin}{width=\maxwidth,height=\maxheight,keepaspectratio}
% Set default figure placement to htbp
\makeatletter
\def\fps@figure{htbp}
\makeatother

\usepackage{booktabs}
\usepackage{longtable}
\usepackage{array}
\usepackage{multirow}
\usepackage{wrapfig}
\usepackage{float}
\usepackage{colortbl}
\usepackage{pdflscape}
\usepackage{tabu}
\usepackage{threeparttable}
\usepackage{threeparttablex}
\usepackage[normalem]{ulem}
\usepackage{makecell}
\usepackage{xcolor}
\usepackage{booktabs}
\usepackage{longtable}
\usepackage[a4paper, left=0.75in, right=1in, top=1in, bottom=1in]{geometry}
\makeatletter
\@ifpackageloaded{caption}{}{\usepackage{caption}}
\AtBeginDocument{%
\ifdefined\contentsname
  \renewcommand*\contentsname{Table of contents}
\else
  \newcommand\contentsname{Table of contents}
\fi
\ifdefined\listfigurename
  \renewcommand*\listfigurename{List of Figures}
\else
  \newcommand\listfigurename{List of Figures}
\fi
\ifdefined\listtablename
  \renewcommand*\listtablename{List of Tables}
\else
  \newcommand\listtablename{List of Tables}
\fi
\ifdefined\figurename
  \renewcommand*\figurename{Figure}
\else
  \newcommand\figurename{Figure}
\fi
\ifdefined\tablename
  \renewcommand*\tablename{Table}
\else
  \newcommand\tablename{Table}
\fi
}
\@ifpackageloaded{float}{}{\usepackage{float}}
\floatstyle{ruled}
\@ifundefined{c@chapter}{\newfloat{codelisting}{h}{lop}}{\newfloat{codelisting}{h}{lop}[chapter]}
\floatname{codelisting}{Listing}
\newcommand*\listoflistings{\listof{codelisting}{List of Listings}}
\makeatother
\makeatletter
\makeatother
\makeatletter
\@ifpackageloaded{caption}{}{\usepackage{caption}}
\@ifpackageloaded{subcaption}{}{\usepackage{subcaption}}
\makeatother

\ifLuaTeX
  \usepackage{selnolig}  % disable illegal ligatures
\fi
\usepackage{bookmark}

\IfFileExists{xurl.sty}{\usepackage{xurl}}{} % add URL line breaks if available
\urlstyle{same} % disable monospaced font for URLs
\hypersetup{
  pdftitle={Beta diversity},
  pdfauthor={Muluh},
  colorlinks=true,
  linkcolor={blue},
  filecolor={Maroon},
  citecolor={Blue},
  urlcolor={Blue},
  pdfcreator={LaTeX via pandoc}}


\title{Beta diversity}
\author{Muluh}
\date{}

\begin{document}
\maketitle


\section{Community composition}\label{community-composition}

\subsection{Phylum}\label{phylum}

\begin{center}
\includegraphics[width=1\textwidth,height=\textheight]{beta_files/figure-pdf/community-composition-phyla-1.png}
\end{center}

\begin{center}
\includegraphics[width=1\textwidth,height=\textheight]{beta_files/figure-pdf/community-composition-phyla-2.png}
\end{center}

\subsection{Genus}\label{genus}

\begin{center}
\includegraphics[width=1\textwidth,height=\textheight]{beta_files/figure-pdf/community-composition-genus-1.png}
\end{center}

\begin{center}
\includegraphics[width=1\textwidth,height=\textheight]{beta_files/figure-pdf/community-composition-genus-2.png}
\end{center}

\subsection{Heatmap}\label{heatmap}

\begin{center}
\includegraphics[width=1\textwidth,height=\textheight]{beta_files/figure-pdf/heatmap-1.png}
\end{center}

\section{Unsupervised analysis}\label{unsupervised-analysis}

\subsection{Community dissimilarity}\label{community-dissimilarity}

This document presents the results of a microbiota project, focusing on
beta diversity analysis. Specifically, in this section, we employed
Principal Coordinates Analysis (PCoA) to assess community dissimilarity
among microbiota samples.

\begin{itemize}
\tightlist
\item
  \textbf{Taxonomic level:} species
\item
  \textbf{Dissimilarity index:} bray
\item
  \textbf{Group variable:} diet
\item
  \textbf{Multiple testing correction:} fdr
\item
  \textbf{Statistical test} Wilcoxon
\end{itemize}

\subsection{PCoA plots}\label{pcoa-plots}

\includegraphics{beta_files/figure-pdf/pcoa-plot-1.png}

\begin{center}
\includegraphics{beta_files/figure-pdf/viz-div_change-1.png}
\end{center}

\subsection{Table of the top findings}\label{table-of-the-top-findings}

\paragraph{Paired treatment between timepoint(1 and
2)}\label{paired-treatment-between-timepoint1-and-2}

\begin{itemize}
\tightlist
\item
  Multiple testing correction: \textbf{paired fdr}
\item
  Statistical test: \textbf{Wilcoxon}
\end{itemize}

\begin{longtable}[t]{lrrr}
\caption{Table of divergence paired between timepoint(1 and 2)}\\
\toprule
\textbf{Meal} & \textbf{mean} & \textbf{p.value} & \textbf{p.adj}\\
\midrule
\endfirsthead
\caption[]{Table of divergence paired between timepoint(1 and 2) \textit{(continued)}}\\
\toprule
\textbf{Meal} & \textbf{mean} & \textbf{p.value} & \textbf{p.adj}\\
\midrule
\endhead

\endfoot
\bottomrule
\endlastfoot
oat & 0.3426224 & 4e-07 & 4e-07\\
rice & 0.3612592 & 3e-07 & 4e-07\\*
\end{longtable}

\subsection{Significant variables}\label{significant-variables}

\includegraphics{beta_files/figure-pdf/significant variables-1.png}

\subsection{Significant taxa}\label{significant-taxa}

\begin{itemize}
\tightlist
\item
  RDA
\end{itemize}

\begin{center}
\includegraphics{beta_files/figure-pdf/significant-taxa-1.png}
\end{center}

\begin{itemize}
\tightlist
\item
  MDS
\end{itemize}

\begin{center}
\includegraphics{beta_files/figure-pdf/mds-1.png}
\end{center}

\section{Supervised analysis}\label{supervised-analysis}

\subsection{Permanova analysis}\label{permanova-analysis}

\begin{itemize}
\item
  Dissimilarity index: \textbf{bray}
\item
  Group variable: \textbf{diet}
\end{itemize}

\textbf{Interpretation}

\begin{itemize}
\item
  Homogeneity \textless{} 0.05: The assumption of homogeneity of
  variance is violated (i.e., the variances are significantly different
  across groups).
\item
  Homogeneity ≥ 0.05: The assumption of homogeneity of variance is not
  violated (i.e., the variances are similar across groups).
\end{itemize}

\subsubsection{Across-group}\label{across-group}

\begin{table}

\caption{PERMANOVA Results - before treatment}
\centering
\begin{tabular}[t]{lrrrr}
\toprule
\textbf{Group1} & \textbf{Group2} & \textbf{P\_value} & \textbf{Homogeneity} & \textbf{p.adj}\\
\midrule
oat & rice & 0.002 & 0.855 & 0.002\\
\bottomrule
\end{tabular}
\end{table}

\paragraph{After treatment}\label{after-treatment}

\begin{table}

\caption{PERMANOVA Results - after treatment}
\centering
\begin{tabular}[t]{lrrrr}
\toprule
\textbf{Group1} & \textbf{Group2} & \textbf{P\_value} & \textbf{Homogeneity} & \textbf{p.adj}\\
\midrule
oat & rice & 0.001 & 0.872 & 0.001\\
\bottomrule
\end{tabular}
\end{table}

\subsubsection{Within-group}\label{within-group}

\begin{table}

\caption{PERMANOVA Results - Paired between timepoint (1 and 2)}
\centering
\begin{tabular}[t]{lrrr}
\toprule
\textbf{Diet} & \textbf{P\_value} & \textbf{Homogeneity} & \textbf{p.adj}\\
\midrule
oat & 1.000 & 0.619 & 1\\
rice & 0.869 & 0.817 & 1\\
\bottomrule
\end{tabular}
\end{table}

\section{Alternate analysis}\label{alternate-analysis}

\includegraphics{beta_files/figure-pdf/viz-fig3-1.png}

\begin{table}[!h]
\centering
\caption{PERMANOVA}
\centering
\begin{tabular}[t]{llrrrrr}
\toprule
\textbf{ } & \textbf{before} & \textbf{after} & \textbf{change-1} & \textbf{change-2} & \textbf{change-3} & \textbf{change-4}\\
\midrule
\cellcolor{gray!10}{betadisp\_pval} & \cellcolor{gray!10}{0.876} & \cellcolor{gray!10}{0.931} & \cellcolor{gray!10}{NA} & \cellcolor{gray!10}{NA} & \cellcolor{gray!10}{NA} & \cellcolor{gray!10}{NA}\\
permanova\_pval & 0.333 & 0.105 & 0.999 & 0.976 & 0.997 & 1\\
\bottomrule
\end{tabular}
\end{table}

\subsection{Data Summary}\label{data-summary}

\begin{verbatim}
class: TreeSummarizedExperiment 
dim: 1700 140 
metadata(0):
assays(2): counts relabundance
rownames(1700): SGB4933 SGB4285 ... SGB15018 SGB4921
rowData names(9): kingdom phylum ... strain clade_name
colnames(140): AE1332 AE2332 ... VK1336 VK2336
colData names(17): sample id ... time_difference time_divergence
reducedDimNames(1): PCoA_BC
mainExpName: NULL
altExpNames(20): kingdom phylum ... KO metacyc
rowLinks: NULL
rowTree: NULL
colLinks: NULL
colTree: NULL
\end{verbatim}




\end{document}
